Lieb and Thirring introduced their eigenvalue inequalities in the study
of fermionic kinetic energy and the stability of matter
\cite{LiebThirring1975,LiebThirring1976}.
The one-state variational problem goes back to Keller \cite{Keller1961}.
In one dimension its optimal constant exceeds the semiclassical constant
for $1/2<\gamma<3/2$. The scalar conjecture in this interval asks whether
summing all negative eigenvalues changes that one-state constant.
The companion paper \cite[Theorem~1.1]{ScalarCompanion} proves that it does
not. The present paper establishes the same constant for arbitrary finite
matrix size, allowing the eigenspaces of the potential to vary with position.

Sharp matrix inequalities at the two endpoint exponents have an
established history. Weidl proved finiteness at the scalar lower endpoint
\cite{Weidl1996}, and Hundertmark, Lieb and Thomas obtained the sharp
$\gamma=1/2$ inequality with constant $1/2$
\cite[Theorem~1]{HLT1998}; Hundertmark, Laptev and Weidl proved the
operator-valued version \cite[Theorem~3.1]{HLW2000}.
At $\gamma=3/2$, Laptev and Weidl established the sharp matrix estimate by
trace identities and extended it to operator-valued potentials
\cite[Theorems~2.1 and~2.2]{LaptevWeidl2000}.
Aizenman--Lieb integration in the spectral parameter gives the larger
exponents \cite{AizenmanLieb1978}. Laptev and Weidl's dimension-lifting
argument then yields the semiclassical constant in every spatial dimension
for $\gamma\ge3/2$ \cite[Theorem~3.1]{LaptevWeidl2000}.
Benguria and Loss gave a proof of the matrix upper-endpoint estimate by
Riccati factorization and commutation \cite{BenguriaLoss2000}.

In the open interval, Hundertmark, Laptev and Weidl obtained the
operator-valued upper bound with twice the semiclassical constant
\cite[Theorem~4.1]{HLW2000}.
Dolbeault, Laptev and Loss developed the method of orthonormal
vector-valued functions at exponent one
\cite[Theorems~1 and~3]{DLL2008}.
More recently, Read and Schulz proved the sharp value
$4/(3\sqrt3\pi)$ for $\gamma=1$, including operator-valued potentials
\cite[Proposition~4 and Equation~(2.1)]{ReadSchulz2026}.
Thus the exponent-one case of our theorem is already contained in their
result. Our additional range is $1/2<\gamma<3/2$, $\gamma\ne1$.

The proof adapts the finite-interval action and continuation framework
of \cite[Theorem~2.1, Proposition~5.1 and Section~6]{ScalarCompanion}.
The companion's comparison applies when the pointwise density has rank
at most one \cite[Proposition~4.1]{ScalarCompanion}.
The new analytic step removes that restriction, making the action method
applicable to vector-valued wave functions.

